\begin{frame}{Table of Contents}
\begin{itemize}
    \item PyTorch CNN Implementation
    \item Reproducibility
    \item Data Handling
    \item Data Augmentation
    \item Transfer Learning
    \item Ensembling
    \item Dropout
    \item Batch Normalization
    \item Full Training Workflow
    \item Evaluation Metrics
\end{itemize}
\end{frame}

\begin{frame}{Learning Outcomes}
\begin{itemize}
    \item Build and train a convolutional network in PyTorch
    \item Make a training run reproducible: seed every generator, control determinism, and checkpoint enough state to resume
    \item Implement efficient data handling techniques using PyTorch's Dataset and DataLoader classes
    \item Apply appropriate data augmentation strategies based on error analysis
    \item Utilize transfer learning effectively for different scenarios and dataset sizes
    \item Understand ensemble methods to improve model performance
    \item Apply regularization techniques like Dropout and Batch Normalization
    \item Design a complete deep learning training workflow from initial setup to inference
    \item Choose and interpret evaluation metrics that suit the task and the class balance
\end{itemize}
\end{frame}

\begin{frame}{Practical Deep Learning}
\begin{itemize}
    \item Practical Implementation of Deep Learning algorithms is just as much an art as it is a science.
    \item The main takeaway is not to start from scratch but rather to build on top of previous knowledge.
    \item Today, we will look at some important tools used in the practical implementation of Deep Learning algorithms.
\end{itemize}
\end{frame} 